\documentclass[leqno]{jsarticle}
\usepackage{amssymb}
\usepackage{amsthm}
\usepackage{amsmath}
\usepackage{comment}
	%可換図式用
		%\usepackage[all]{xy}
	%重くなるので使わいときはコメント化しておく。
%定理環境 thmはあまり使わずpropをなるべく使う。
%主張は基本的に証明の中で使い、そのため番号を付けない。
%注意環境を付けてみた。
\newtheorem{thm}{定理}
\newtheorem{prop}[thm]{命題}
\newtheorem{cor}[thm]{系}
\newtheorem{lem}[thm]{補題}
\newtheorem{df}[thm]{定義}
\newtheorem{exam}[thm]{例}
\newtheorem{fact}[thm]{事実}
\newtheorem*{claim}{主張}
\newtheorem*{cau}{注意}
\renewcommand{\proofname}{{\bf 証明}}
%矢印たち。
%\toは\rightarrowと同じ。\getsは\leftarrowと同じ。同値は\iff
%上向き矢はuparrow逆はdownarrow
\newcommand{\To}{\Longrightarrow}
\newcommand{\Gets}{\Longleftarrow}
%黒板太字
\newcommand{\nn}{\mathbb{N}}
\newcommand{\zz}{\mathbb{Z}}
\newcommand{\qq}{\mathbb{Q}}
\newcommand{\rr}{\mathbb{R}}
\newcommand{\cc}{\mathbb{C}}
%無限大
\newcommand{\mgn}{\infty}
%ギリシャン
\newcommand{\dpst}{\displaystyle}
\newcommand{\ep}{\varepsilon}
\newcommand{\al}{\alpha}
\newcommand{\be}{\beta}
\newcommand{\de}{\delta}
\newcommand{\la}{\lambda}
\newcommand{\La}{\Lambda}
\newcommand{\ga}{\gamma}
\newcommand{\lil}{\la\in \La}
%商集合
\newcommand{\qt}[2]{#1/\! #2}
%enumerate環境
\renewcommand{\labelenumi}{{\rm (\arabic{enumi})}}
%以降alignじゃなくてenumerate を使え。
%なんか演算
\newcommand{\card}{\text{  {\rm card}} }
\newcommand{\supp}{\text{  {\rm supp}} }
%なんか演算２
\newcommand{\bd}{\text{{\rm Bdry}} }
\newcommand{\cl}{\text{{\rm CL}} }
\newcommand{\nai}{\text{{\rm Int}} }
%%%%%%%%%%%%%%%%%%%%%%%%%%%%%%%%%%%%%
\newcommand{\s}{\text{  {\rm St}} }
\newcommand{\uu}{\mathcal{U}}
\newcommand{\vv}{\mathcal{V}}
\newcommand{\A}{\mathcal{A}}
\newcommand{\B}{\mathcal{B}}
\newcommand{\C}{\mathcal{C}}
\newcommand{\D}{\mathcal{D}}
\newcommand{\E}{\mathcal{E}}
\newcommand{\W}{\mathcal{W}}
\newcommand{\dd}{\Delta}
\newcommand{\pr}{\prec}
\newcommand{\zp}{\zz_{+}}
\newcommand{\zr}{\zz_{\ge0}}

%差集合は\setminusで統一する。等号の否定は\neq
%真部分集合は\subsetneqq　偏微分は\partial
%ディスプレイスタイル。\dfracでディスプレイ分数
%空集合は\Oを使う！！！！！！！！！！！！

\title{距離化可能定理：Alexandroff-Urysohn-Tukey}
\author{一色三良(concious77)}
\date{\today}
			\begin{document}
\maketitle

この文章では所謂Alexandroff-Urysohn-Tukeyの距離化可能定理を証明する。そのためにまず最初のセクションでは被覆についての基本的な事項を紹介する。

\section*{正規被覆}

\begin{df}[星型集合]
集合$X$とその部分集合族$\mathcal{A}$、及び$S\subset X$について
\[
\s(S,\mathcal{A})=\bigcup\{A\ |\ A\in \mathcal{A},\ S\cap A\neq\O\}
\]
を$\mathcal{A}$による$S$を中心とする星型集合という。
\end{df}

\begin{df}[$*$と$\Delta$]
$X$の被覆$\uu$について$\uu^{*}$を
\[
\uu^{*}=\{\s(S,\uu)\ |\ S\in \uu\}
\]
と定義し、そして
$\uu^{\Delta}$を
\[
\uu^{\Delta}=\{\s(x,\uu)\ |\ x\in X\}
\]
と定義する。
\end{df}


\begin{df}[細分]
$\mathcal{A},\mathcal{B}を$集合$X$の部分集合の族とする。このとき$\mathcal{A}$が$\mathcal{B}$の細分であるとは
\[
\forall B\in \mathcal{B}\ \exists A\in \mathcal{A}\ A\subset B
\]
となることである。\par
$\A$が$\B$の細分であるとき
\[
\A\pr \B
\]
と書くことにしよう。
\end{df}


\begin{df}[星型細分、$\Delta$細分]
集合$X$の部分集合族$\mathcal{A},\B$について$\A$が$\B$の星型細分であるとは
\[
\A^{*}\pr \B
\]
となることである。\par
また、集合$X$の部分集合族$\mathcal{A},\B$について$\mathcal{A}$が$\B$の$\Delta$細分、又は重心細分であるとは
\[
\A^{\dd}\pr\B
\]
\end{df}

\begin{lem}\label{benri}
$A\subset B,\uu\pr \mathcal{V}$ならば
\[
\s(A,\uu)\subset \s(B,\mathcal{V})
\]
\end{lem}
\begin{proof}
まず$\s(A,\uu)\subset \s(B,\uu)$は明らかである。次に$\s(B,\uu)\subset \s(B,\vv)$を示そうとするのだが、細分なので明らかである。よって証明は終わる。
\end{proof}


\begin{lem}\label{aaa}$X$の被覆$\A,\B$について
\[
\A^{\dd}\pr \B\To \A\pr\B
\]
\end{lem}
\begin{proof}
自明
\end{proof}




\begin{lem}\label{aaa}$X$の被覆$\A,\B$について
\[
\A^{*}\pr \B\To \A^{\dd}\pr\B
\]
\end{lem}
\begin{proof}
補題\ref{benri}から分かる。
\end{proof}

\begin{lem}\label{bebe}$X$の被覆$\A,\B,\C$
\[
\A^{\dd}\pr\B,\B^{\dd}\pr\C\To \A^{*}\pr\C
\]
\end{lem}
\begin{proof}
$A\in \A$を任意に与える。$a\in A$を適当に取る。$V\in \A,A\cap V\neq\O$とすると$\A^{\dd}\pr\B$よりある$B\in\B$が存在し、$A\cup V\subset B$となる。このとき$a\in B$でもある。よってこれから
\[
\s(A,\A)\subset \s(a,\B)
\]
そして$\B^{\dd}\pr \C$なので$\A^{*}\pr \C$となることはすぐ分かる。
\end{proof}

\begin{lem}
$X$の部分集合$S$と$X$の被覆$\A,\B$について
$\A^{*}\pr \B$のとき
\[
\s(\s(S,\A),\A)\subset  \s(S,\B)
\]
が成り立つ。
\end{lem}
\begin{proof}
明らか
\end{proof}


\begin{df}[正規列、正規被覆]
$X$の開被覆の列$\{\uu_n\}_{n\in \zr}$が正規列であるとは任意の$n$について
\[
\uu_{n+1}^{\dd}\pr \uu_n
\]
となること。\par
また、開被覆$\mathcal{V}$が正規被覆であるとは正規列$\{\uu_n\}_{n\in \zr}$が存在して
\[
\uu_0\pr\mathcal{V}
\]
となること。
\end{df}


\setcounter{equation}{0}
\begin{lem}\label{ccc}
$X$の被覆$\mathcal{V}$が正規被覆であることと、被覆の列$\{\uu_n\}_{n\in \zr}$が存在して
\[
\uu^{*}_{n+1}\pr\uu_n,\ \uu_0\pr\mathcal{V}
\]
を充たすことは同値である。

\end{lem}
\begin{proof}
補題\ref{aaa},\ref{bebe}からすぐに分かる。
\end{proof}


\begin{lem}\label{hey}
位相空間$X$の開被覆$\uu$と集合$A$と開集合$U$について
\[
\s(A,\uu)\subset U
\]
が成り立つならば$\cl(A)\subset U$となる。
\end{lem}
\begin{proof}
$a\in \cl(A)$を任意に与えると$\uu$が被覆であることから$V\in \A,a\in V$となる$V$が存在するが、$a$は$A$の触点なので$V\cap A\neq \O$である。よって$V\subset \s(A,\uu)\subset U$である。故に$a\in U$である。以上から$\cl(A)\subset U$が分かる。
\end{proof}

\setcounter{equation}{0}
\begin{thm}[Stone]\label{stone}
位相空間$X$の正規被覆には局所有限な開細分被覆が存在する。
\end{thm}
\begin{proof}
%%%%%%%%%%%%%
$\uu$を正規被覆とした時、補題\ref{ccc}より星型細分の列が存在するが、適当に番号を付け替えることによって開被覆$\vv$と開被覆列$\{\A_n\}_{n\in \zr}$が存在して
\[
\A_0\pr \vv,\ \vv^{*}\pr\uu
\]
でかつ
\[
\A_{n+1}^{*}\pr\A_n
\]
を充たす。\par
そして帰納的に
\[
\B_1=\A_1,\ \B_2=\{\s(V,\A_1)\ |\ V\in \B_1\},\dots,\B_n=\{\s(V,\A_n)\ |\ V\in \B_{n-1}\},\dots
\]
と置く。この時各$n\in \zp$について$\{\s(V,\A_n)\ |\ V\in\B_n\}$は$\A_0$の細分である。このことを帰納法によって証明しよう。まず$n=1$の時は明らかである。次に$n=k$までは成り立っているとして$n=k+1$のとき、$V\in \B_{k+1}$を与え$U=\s(V,\A_{k+1})$を考える。この時$V$は$\B_{k+1}$の作り方から$S\in \B_k$を用いて$V=\s(S,\A_{k+1})$と表される。そして$\A_{k+1}^{*}\pr \A_k$なので
\[
U=\s(V,\A_{k+1})=\s(\s(S,\A_{k+1}),\A_{k+1})\subset \s(S,\A_k)
\]
なので帰納法の仮定から$k+1$の場合も成り立つ。\par
以上のことから特に各$\B_n$は$\A_0$の細分である。\par
%%第二弾楽
次に$X$を整列させてその順序を$<$とする。そして$(n,x)\in \zp\times X$について
\[
C_n(x)=\s(x,\B_n)\setminus \bigcup_{y<x}\s(y,\B_{n+1})
\]
と定義する。この時$\C=\{C_n(x)\ |\ n\in \zp,x\in X\}$は$X$の被覆となる。実際$p\in X$について
\[
P=\{z\in X\ |\ p\in \bigcup_{n\in \zp}\s(z,\B_n)\}
\]
とを考えると$p\in P$なので$P$は空では無いので最小限$y$が存在する。$y$についてある$n$が存在して$p\in \s(y,\B_n)$となる。このような$n$について$p\in C_n(y)$となる。よって$\C$は被覆となる。\par
次の主張が成り立つ。
%%%%%%%%%%%%%%%%%%%%%%%%%%
\begin{claim}[a]
$n\in \zp$を固定したとき、$U\in \A_{n+1}$は$\{C_n(x)\ |\ x\in X\}$の高々１個の元としか交わらない
\end{claim}
{\bf [主張(a)の証明]}\par
いま、$U\cap C_n(x)\neq\O$としよう。すると$x\in V$となる$V\in \B_n$が存在して$U\cap V\neq\O$よって
\[
x\in U\cup V\subset \s(V,\A_{n+1})\in \B_{n+1}
\]
故に$U\subset \s(x,\B_{n+1})$よって$\C_n(x)$の作り方から$x$を$U\cap C_n(x)\neq\O$となる最小の元とすれば$U$は$C_n(x)$以外の$\{C_n(x)\ |\ x\in X\}$の元と交わらない。\par
{\bf [主張(a)の証明終わり]}\par
%%%%%%%%%%%%%%%%%%%%%%%%%
また、$C_n(x)\subset \s(x,\B_{n})$であり、$\B_n\pr \A_0$でかつ$\A_0^{*}\pr \vv$なので$\C\pr\vv$である。\footnote{証明が終わったあと詳しく説明する}\par
%%%%%%%%%%%%%%%%%%%%%%%%%%
さて$(n,x)\in \zp\times X$について
\[
D_n(x)=\s(C_n(x),\A_{n+3}),\ E_n(x)=\s(C_n(x),\A_{n+2})
\]
と定義し、
\[
\D=\{D_n(x)\ |\ (n,x)\in \zp\times X\},\ \E=\{E_n(x)\ |\ (n,x)\in \zp\times X\}
\]
と定義する。\par
このとき
\[
\s(D_n(x),\A_{u+3})=\s(\s(C_n(x),\A_{n+3}),\A_{n+3})\subset \s(C_n(x),A_{n+2})=E_n(x)
\]
が成り立つので補題\ref{hey}から$\cl(D_n(x))\subset E_n(x)$である。\par
ここで次の主張が成り立つ。
%%%%%%%%%%%%%%%%%%%%%%%%%%%%%%%%
\begin{claim}[b]
$n\in \zp$を固定すると$U\in \A_{n+2}$は$\{E_n(x)\ |\ x\in X\}$の高々１個の元とし交わらない。特に$\{E_n(x)\ |\ x\in X\}$は疎である。
\end{claim}
{\bf [主張(b)の証明]}\par
$U\cap E_n(x)\neq\O$とすると$E_n(x)$の定義から$V\cap C_n(x)\neq\O,\ U\cap V\neq\O$となる$V\in \A_{n+2}$が存在するが、$\uu_{n+2}^{*}\pr \uu_{n+1}$なので$U\cup V\subset O$となる$O\in \A_{n+1}$が存在し、主張$(a)$から$O$は$C_n(x)$以外の$\{C_n(x)\ |\ x\in X\}$の元と交わらないので、$U$は$\{E_n(x)\ |\ x\in X\}$の高々１個の元としか交わらない。\par
{\bf [主張(b)の証明終わり]}\par
%%%%%%%%%%%%%%%%%%%%%%%%%%%%%
また、$D_n(x)\subset E_n(x)$なので主張(b)から$n$を固定すると$\{D_n(x)\ |\ x\in X\}$は疎である。よって
\[
F_n=\bigcup_{x\in X}\cl(E_n(x))
\]
は閉集合である。\par
%%%%%%%%%%%%%%%%%%%%%%%%%%%%%
そして$\C\pr \vv,\ \A_n\pr \vv$で$\vv^{*}\pr\uu$なので$\mathcal{E}\pr \uu$である。\footnote{証明が終わった後詳しく説明する。}\par
%%%%%%%%%%%%%%%%%%%%%%%%%
以上を元に$(n,x)\in \zp\times X$について
\[
W_n(x)=E_n(x)\setminus \bigcup_{i<n}F_i
\]
と定義する。$\mathcal{W}=\{W_n(x)\ |\ (n,x)\in \zp\times X\}$が$\uu$の局所有限な開細分被覆であることを示そう。\par
まず$\mathcal{E}\pr \uu$であり、$W_n(x)\subset E_n(x)$なので$\mathcal{W}\pr \uu$である。\par
次に$\mathcal{W}$が被覆であることを示そう。$p\in X$を任意に与える。
\[
\exists x\in X\ p\in \cl(D_n(x))
\]
を成り立たせる最小の$n$を$m$とすると$y\in X$が存在して$p\in \cl(D_m(y))$となり、$m$の最小性から\\
$p\not\in F_1\cup F_2\cup\cdots \cup F_{m-1}$で、かつ$p\in E_m(y)$となる。よって$p\in W_m(y)$なので$\mathcal{W}$は$X$の被覆となる。\par
最後に$\mathcal{W}$の局所有限性を示そう。$p\in X$を任意に与える。すると$p\in C_i(y)$となる$C_i(y)\in \C$が存在する。このとき
\[
\s(p,\A_{i+3})\subset D_i(y)
\]
なので$\s(p,\A_{i+3})\subset F_i$よって$k>i$ならば$\s(p,\A_{i+3})$は$W_k(x)\ (x\in X)$と交わらない。\par
そして$k< i$の時には$\A_i^{*}\pr \A_k$であり、特に$\A_k^{\dd}\pr \A_i$
なので$\s(p,\A_{i+3})\subset V_{k}$となる
$V_k\in \A_{k}$が存在する。以上と主張$(b)$を踏まえると$p\in N$となる$N\in \A_{i+3}$を適当に取れば、$N\subset V_k$なので
$N$は$\mathcal{W}$の高々$i$個の元としか交わらない。\par
これで証明が終わった。
\end{proof}
\begin{cau}[証明の中で後で説明すると言った箇所]${}$\\
{\bf [$C_n(x)\subset \s(x,\B_{n})$であり、$\B_n\pr \A_0$で$\A_0^{*}\pr \vv$なので$\C\pr\vv$である。]}の説明：$x\in M$となる$M\in \A_0$を取ると補題\ref{benri}から
\[
C_n(x)\subset \s(x,\B_{n})\subset \s(M,\A_n)\in \A_0^{*}
\]
$\A_0^{*}\pr\vv$なので$\s(M,\A_n)\subset O$となる$O\in \vv$が存在するので結局
\[
C_n(x)\subset O
\]
となる$O\in \vv$が存在することになる。よって$\C\pr \vv$


{\bf [$\C\pr \vv,\ \A_n\pr \vv$で$\vv^{*}\pr\uu$なので$\mathcal{E}\pr \uu$である。]}の説明：
$E_n(x)=\s(C_n(x),\A_{n+2})$である。$\C\pr \vv$から$C_n(x)\subset V$となる$V\in \vv$が存在する。これと$\A_{n+2}\pr \vv$と補題\ref{benri}から
\[
E_n(x)=\s(C_n(x),\A_{n+2})\subset \s(V,\vv)\in \vv^{*}
\]
である。また、$\vv^{*}\pr \uu$なので$\s(V,\vv)\subset O$となる$O\in \uu$が存在するので結局
\[
E_n(x)\subset O
\]
となる$O\in \uu$が存在することになる。よって$\E\pr \uu$
\end{cau}



\section*{距離化可能定理：被覆}
\begin{fact}[Bing-長田-Smirnovの距離化可能定理]\label{BNS}
$T_3$空間$X$について以下は同値
\begin{align}
&\mbox{$X$は距離化可能}\\
&\mbox{$X$は$\sigma$疎な開基を持つ}\\
&\mbox{$X$は$\sigma$局所有限な開基を持つ}
\end{align}
\end{fact}



\begin{thm}[Alexandroff-Urysohn-Tukeyその0]
$T_1$空間$X$について以下の条件は距離化可能性と同値。\par
$X$の被覆の列$\{\mathcal{U}_n\}_{n\in \nn}$が存在し、各$n\in \nn$について$\mathcal{U}_n$は正規被覆であり、そして各$x\in X$について
$\{S(x,\mathcal{U}_n)\}_{n\in \nn}$が$x$の基本近傍系になる。
\end{thm}
\begin{proof}
距離化可能性からこの条件が従うのは$2^{-n}$近傍とかを考えれば分かる。逆を示そう。\par
まず$X$が$T_3$である事を示そう。$C$を閉集合として$p\not\in C$とし、$n$を
\[
\s(p,\uu_n)\subset X\setminus C
\]
を成り立たせるものとする。この時$\uu_n$が正規被覆であることから$\A^{\dd}\pr\uu_n$となる開被覆$\A$が存在する。そして
\[
P=\s(p,\A),\ Q=\s(C,\A)
\]
と置くと$P\cap Q=\O$である。実際、$P\cap Q\neq \O$と仮定すると$p\in N,\ C\cap M\neq \O,\ N\cap M\neq\O$となる$N,M\in \A$が存在するが、$x\in N\cap M$について$\A^{\dd}\pr \uu_n$より$\s(x,\A)\subset O$となる$O\in\uu_n$が存在し、$p\in O\subset \s(p,\uu_n)$で$O\cap C\neq\O$となるが、$\s(p,\uu_n)\subset X\setminus C$に反する。以上から$X$は$T_3$空間である。\par
さて距離化可能性を示そう。各$\uu_n$について定理\ref{stone}から$\vv_n\pr\uu_n$となる局所有限開被覆$\vv_n$が存在する。このとき$\{\s(x,\vv_n)\}_{n\in \nn}$は$x$の基本近傍系になる。よって$U$を任意に与えると点$p\in U$についてある$k$が存在して$\s(p,\vv_k)\subset U$となる。故に$p\in V\in U$となる$V\in \vv_n$が存在する。この事から$\bigcup_{n\in \nn}\vv_n$は$X$の開基になる。そして事実\ref{BNS}から$X$は距離化可能である。
\end{proof}



\begin{thm}[Alexandroff-Urysohn-Tukeyその１]
$T_1$空間$X$について以下の条件は距離化可能性と同値。\par
$X$の正規列$\{\mathcal{U}_n\}_{n\in \nn}$が存在し、各$x\in X$について
$\{S(x,\mathcal{U}_n)\}_{n\in \nn}$が$x$の基本近傍系になる。
\end{thm}
\begin{proof}
正規列$\{\uu_n\}_{n\in\nn}$に属する各々の$\uu_m$は正規被覆になるので先の定理から定理は明らかである。
\end{proof}














\begin{thebibliography}{9}
\bibitem{1}森田紀一,位相空間論,岩波全書331,岩波書店1981
\bibitem{d}J.Dugundji,{\it Topology},William C Brown Pub ,1966
\bibitem{stone}A.H.Stone,{\it Paracompactness and product spaces},Bull,Amer.Math Soc. 54(1948),977-982
\bibitem{Tukey}J.W.Tukey,{\it Convergence and uniformity in topology},Annals of Mathematics Studies, \par no.2,Princeton,1940

\end{thebibliography}










			\end{document}



\begin{comment}
[可換図式]
\[\xymatrix{
&   & (2) \ar@{=>}[rd]&  &  &  & (5) \ar@{=>}[rd]&\\ 
& (1)\ar@{=>}[ru] \ar@{=>}[rd]&  &(4)\ar@{=>}[r]&(7)& (1)\ar@{=>}[ru] \ar@{=>}[rd]& & (7)&\\ 
&   & (3) \ar@{=>}[ur]&  &  &  &(6) \ar@{=>}[ur]&
}\]

[\bigin \endを使わない方法]

{\prop[aa] aaa}
で
\begin{prop}[aa]
aaa
\end{prop}
と同じ\df \thm \proof でも同様

[直和]
$\amalg$

[同値関係でよく使う波線]
\sim

[引数付きの新命令]
\newcommand{\prn}[1]{\|#1\|_p}

[番号のリセット]

\setcounter{equation}{0}


[字体の変更]

{\rm hogehoge}と使う。

\rm ローマン
\it イタリック
\sl 斜体

[supp]

\newcommand{\supp}{\text{{\rm supp}}}

[中央揃えの改行]

	\begin{eqnarray*}
		hoge1&=&hogehoge
		hoge2&=&hogehofe

	\end{eqnarray*}

&&の部分で揃えられる。






[場合分け]

	\begin{cases}
		hoge1 & \text{状況１}\\
		hoge2 & \text{状況２}\\
		・
		・
		・
	\end{cases}



\end{comment}





